
Reward models guide RLHF training, yet aggregate evaluation metrics conceal systematic failure patterns. This paper introduces a $2\times 2$ mode decomposition that classifies preference battles by human vote entropy (high/low) crossed with reward model variance (high/low), yielding four alignment modes. Analysis of 57,477 Chatbot Arena battles reveals that 23.7\% fall into Mode 3 (Misaligned-Confident): samples where human voters disagree while the reward model expresses high confidence (95\% CI: 23.4--24.0\%, p < 0.001 against a 10\% threshold). Two mechanistic hypotheses were tested: (1) that Mode 3 arises from semantic divergence between responses, and (2) that Mode 3 concentrates in subjective prompt types. Both hypotheses are unsupported: Mode 3 pairs exhibit higher semantic similarity than aligned pairs (Cohen's d = $-0.049$, opposite to prediction), and Mode 3 rates are nearly identical across subjective and objective tasks (ratio = 1.07, below the 1.5 threshold). The contributions are: (a) the first empirical quantification of overconfident reward model misalignment in large-scale preference data, (b) a mode decomposition framework as an alignment diagnostic, and (c) negative results ruling out two candidate mechanisms.
